%---------------------------------------------------------------------
% Quick-reference sheet and the recurring mistakes.
%---------------------------------------------------------------------

\section*{فوری دہرائی}
\markboth{فوری دہرائی}{}

\subsection*{خبر کا ڈھانچہ --- ایک نظر میں}

\begin{center}
\begin{tabular}{@{}p{3.2cm} p{11.2cm}@{}}
\toprule
سرخی & خلاصہ چند الفاظ میں --- سب ایڈیٹر لگاتا ہے \\
ڈیٹ لائن & جگہ اور تاریخ: ''لاہور، \textenglish{5} ستمبر'' \\
انٹرو (سرخبر) & پہلا پیرا، \textenglish{25--35} الفاظ، سب سے اہم بات \\
باڈی & تفصیل، بیانات، پس منظر --- گھٹتی اہمیت کے ساتھ \\
\midrule
چھ سوالات & کیا، کون، کب، کہاں، کیوں، کیسے \\
الٹا اہرام & اہم بات اوپر؛ خبر نیچے سے کاٹی جا سکے \\
\bottomrule
\end{tabular}
\end{center}

\subsection*{تین اہم فرق جو اکثر پوچھے جاتے ہیں}

\begin{center}
\begin{tabular}{@{}p{3.4cm} p{11.0cm}@{}}
\toprule
\textbf{رپورٹر / نامہ نگار / نمائندہ}
  & مرکزی شہر میں / دوسرے شہر یا ملک میں / مقامی سطح پر \\
\textbf{خبر / تجزیہ / اداریہ}
  & کیا ہوا / کیوں ہوا اور نتائج / ہمیں کیا کرنا چاہیے \\
\textbf{تحقیقی / تجزیاتی}
  & چھپی حقیقت بے نقاب کرنا / کھلے واقعے کا مطلب سمجھانا \\
\bottomrule
\end{tabular}
\end{center}

\subsection*{خبر رساں ادارے --- سنین}

\begin{center}
\begin{tabular}{@{}p{3.6cm} p{2.0cm} p{8.6cm}@{}}
\toprule
\textenglish{Havas} & \textenglish{1835} & فرانس --- دنیا کا پہلا؛ بعد میں \textenglish{AFP} \\
\textenglish{AP} & \textenglish{1846} & امریکہ \\
\textenglish{Wolff} & \textenglish{1849} & جرمنی \\
\textenglish{Reuters} & \textenglish{1851} & برطانیہ --- پال جولیس رائٹر \\
\textenglish{UPI} & \textenglish{1907} & امریکہ \\
\textenglish{TASS} & \textenglish{1925} & روس \\
\textenglish{Xinhua} & \textenglish{1931} & چین \\
\textenglish{APP} & \textenglish{1947} & پاکستان --- سرکاری \\
\bottomrule
\end{tabular}
\end{center}

\medskip
\begin{nukta}[خبر کی نو اقدار --- یاد رکھنے کی فہرست]
تازگی \ \textbullet\ \ قربت \ \textbullet\ \ اہمیت و اثر \ \textbullet\ \
نمایاں شخصیت \ \textbullet\ \ تصادم \ \textbullet\ \ غیر معمولی پن \
\textbullet\ \ انسانی دلچسپی \ \textbullet\ \ نتائج \ \textbullet\ \ تعداد و
حجم۔

\smallskip
کسی بھی خبر کے سوال میں ان میں سے چار پانچ کا ذکر جواب کو مکمل کر دیتا ہے۔
\end{nukta}

\clearpage

%---------------------------------------------------------------------
\section*{آٹھ عام غلطیاں}
\markboth{آٹھ عام غلطیاں}{}

\begin{enumerate}\setlength{\itemsep}{5pt}
  \item \textbf{خبر میں اپنی رائے لکھنا۔} خبر صرف \emph{کیا ہوا} بتاتی ہے۔
        رائے ہو تو کسی ذمہ دار شخص کے حوالے سے --- ''ترجمان کے مطابق''۔
  \item \textbf{رپورٹر اور سب ایڈیٹر کے فرائض ملا دینا۔} رپورٹر خبر
        \emph{لاتا} ہے، سب ایڈیٹر اسے \emph{سنوارتا} ہے اور سرخی لگاتا ہے۔
  \item \textbf{رپورٹر، نامہ نگار اور نمائندہ کو ایک سمجھنا۔} فرق کی بنیاد
        \emph{مقام اور حیثیت} ہے --- مرکز، دوسرا شہر، مقامی سطح۔
  \item \textbf{ملزم کو مجرم لکھ دینا۔} عدالت کے فیصلے سے پہلے ہر شخص
        \emph{ملزم} ہے۔ یہ ہتکِ عزت کا سب سے عام سبب ہے۔
  \item \textbf{تجزیے اور اداریے کو ایک سمجھنا۔} تجزیہ نتائج بتاتا ہے، اداریہ
        رائے دیتا ہے اور مشورہ دیتا ہے۔
  \item \textbf{اقدار اور اقسام میں خلط ملط۔} \emph{اقدار} خبر کی کسوٹی ہیں
        (تازگی، قربت)؛ \emph{اقسام} موضوع کے اعتبار سے تقسیم ہیں (سیاسی،
        کھیل)۔ سوال غور سے پڑھیے۔
  \item \textbf{سنین یاد نہ ہونا۔} خبر رساں اداروں کے سوال میں
        \textenglish{1835}، \textenglish{1846}، \textenglish{1851} اور
        \textenglish{1947} ہی امتیاز پیدا کرتے ہیں۔
  \item \textbf{وقت کی بدانتظامی۔} پانچ سوال، تین گھنٹے --- فی سوال تقریباً
        پینتیس منٹ۔ پہلے دو سوالوں پر ڈیڑھ گھنٹہ لگا دینا سب سے بڑا نقصان ہے۔
\end{enumerate}

\vfill
\begin{center}
\begin{tcolorbox}[enhanced, colback=bmtint, colframe=bmblue, boxrule=0.5pt,
                  arc=3pt, width=0.94\linewidth,
                  left=11pt, right=11pt, top=9pt, bottom=9pt]
\small
\textbf{\color{bmblue}اگر وقت بہت کم ہو۔}\ \
دونوں پرچوں میں دہرائے گئے پانچ موضوعات پہلے تیار کیجیے:
\textenglish{(1)} رپورٹر کے فرائض و اوصاف اور رپورٹر/نامہ نگار/نمائندہ کا فرق
(سوال \textenglish{6}، \textenglish{7})؛
\textenglish{(2)} خبر کا مفہوم اور کردار (سوال \textenglish{1})؛
\textenglish{(3)} مکتوب نگاری (سوال \textenglish{12})؛
\textenglish{(4)} خبر رساں ادارے (سوال \textenglish{13})؛
\textenglish{(5)} ترقی یافتہ و تجزیاتی خبر نگاری (سوال \textenglish{10}،
\textenglish{11})۔
پانچ میں سے تین سوال عموماً انہی سے بن جاتے ہیں۔

\medskip
\textbf{\color{bmblue}ماخذ۔}\ \
اس کورس کی سرکاری درسی کتاب دستیاب نہ ہو سکی، اس لیے تمام سوالات اے آئی او یو
کے دو حقیقی پرچوں سے لیے گئے ہیں: بی اے خبر نگاری (\textenglish{431})، سمسٹر
بہار \textenglish{2013} اور سمسٹر خزاں \textenglish{2013}۔ جوابات خبر نگاری کے
مسلمہ اصولوں کے مطابق ہیں --- کتاب کی مخصوص تفصیلات اپنی درسی کتاب سے ملا
لیجیے۔
\end{tcolorbox}
\end{center}
